\input{style/sqtbeamer.tex}
\usepackage{tikz}

% Dice macros (\die{n}, \diehidden) for the trading sims.
% Reusable six-sided die macro.
% Usage: \die{n}      draws a d6 showing face n (1-6)
%        \die[<opts>]{n}  passes extra tikzpicture options (e.g. scale=0.8)
%
% Pips are laid out on a 2x2 grid inside a rounded square.
\newcommand{\diepip}[2]{\fill[black] (#1,#2) circle (0.16);}

\NewDocumentCommand{\die}{O{} m}{%
  \begin{tikzpicture}[scale=0.7, #1]
    \draw[rounded corners=4pt, fill=white, line width=1pt] (0,0) rectangle (2,2);
    \ifcase#2\relax
      % 0: blank
    \or % 1
      \diepip{1.0}{1.0}
    \or % 2
      \diepip{0.5}{1.5}\diepip{1.5}{0.5}
    \or % 3
      \diepip{0.5}{1.5}\diepip{1.0}{1.0}\diepip{1.5}{0.5}
    \or % 4
      \diepip{0.5}{1.5}\diepip{1.5}{1.5}\diepip{0.5}{0.5}\diepip{1.5}{0.5}
    \or % 5
      \diepip{0.5}{1.5}\diepip{1.5}{1.5}\diepip{1.0}{1.0}\diepip{0.5}{0.5}\diepip{1.5}{0.5}
    \or % 6
      \diepip{0.5}{1.5}\diepip{1.5}{1.5}\diepip{0.5}{1.0}\diepip{1.5}{1.0}\diepip{0.5}{0.5}\diepip{1.5}{0.5}
    \fi
  \end{tikzpicture}%
}

% A blank "unknown" die (question mark face), for hidden rolls.
\NewDocumentCommand{\diehidden}{O{}}{%
  \begin{tikzpicture}[scale=0.7, #1]
    \draw[rounded corners=4pt, fill=black!5, line width=1pt] (0,0) rectangle (2,2);
    \node at (1,1) {\Large\bfseries ?};
  \end{tikzpicture}%
}


\title{Trading Workshop}
\subtitle{Mock Trading with a Live Engine}
\institute{Sydney Quant Traders}
\date{}

\begin{document}

\maketitle

% =====================================================================
% NOTE ON FORMAT
% ---------------------------------------------------------------------
% This deck is a *skeleton*. The slides provide visuals to guide the
% workshop; the actual trading happens in the mock trading engine.
%
% Slides marked "TODO" are intentionally left empty to be filled in.
% Each trading sim ("Trading Sim") is a round played on the engine ---
% the slide is just the prompt / visual shown to the room.
% =====================================================================


% #####################################################################
\section{Introduction to Trading}
% #####################################################################
% Short intro (a condensed version of Intro to Quant). Leave empty for now.

\begin{frame}{What is Trading?}
  % TODO: fill in --- high level: exchanging assets, why we trade.
\end{frame}

\begin{frame}{Buying and Selling}
  % TODO: fill in --- bids/asks, what it means to buy vs sell, price.
\end{frame}

\begin{frame}{Position and Max Position}
  % TODO: fill in --- long/short position, PnL from position,
  % the max position limit enforced by the engine.
\end{frame}

\begin{frame}{How the Trading Engine Works}
  % TODO: fill in --- how players connect, send bids/asks, see the book,
  % how positions are liquidated at the true value at the end of a round.
\end{frame}


% #####################################################################
\section{Trading Sim 1: Sum of Five Dice}
% #####################################################################

\begin{frame}{Trading Sim 1: Sum of Five Dice}
  \begin{sqt:definition}
    The true value of the asset is the \textbf{sum of five dice rolls}.
    $$V = d_1 + d_2 + d_3 + d_4 + d_5, \qquad d_i \sim \mathrm{Unif}\{1,\dots,6\}.$$
  \end{sqt:definition}

  \begin{itemize}
    \item Make a market on $V$ in the trading engine.
    \item Everyone can see the dice --- there is no hidden information yet.
    \item At the end, positions are liquidated at the true $V$.
  \end{itemize}

  \vspace{0.3cm}
  \textcolor{sqt-question}{What is a fair value for $V$? How wide should your market be?}
\end{frame}

\begin{frame}{The Five Dice}
  \begin{center}
    \vspace{0.4cm}
    % TODO: set the rolled values here (currently 3, 5, 2, 6, 4).
    \die{3}\hfill\die{5}\hfill\die{2}\hfill\die{6}\hfill\die{4}

    \vspace{0.8cm}
    {\Large $3 + 5 + 2 + 6 + 4 = \mathbf{20}$}
  \end{center}
\end{frame}

\begin{frame}{Fair Value}
  % TODO: fill in --- E[one die] = 3.5, so E[V] = 5 * 3.5 = 17.5.
  % Discuss variance / how tight a market you can make.
\end{frame}


% #####################################################################
\section{Trading Sim 2: Estimation with Reveals}
% #####################################################################
% An estimation market where information is revealed over rounds.

\begin{frame}{Trading Sim 2: Estimation Game}
  \begin{sqt:definition}
    The true value is some real-world quantity that nobody knows exactly.
    Clues are \textbf{revealed round by round}; update your market as you learn more.
  \end{sqt:definition}

  % TODO: pick the quantity, e.g. "number of X", "volume of Y".
  \textbf{The asset:} \textcolor{sqt-question}{TODO --- state the quantity to estimate.}
\end{frame}

\begin{frame}{Estimation Game: Reveals}
  \textbf{Round 1}
  \begin{itemize}
    \item % TODO: first clue. Players open a market and trade.
  \end{itemize}

  \onslide<2->{
    \textbf{Round 2}
    \begin{itemize}
      \item % TODO: second clue. Players refine and trade.
    \end{itemize}
  }

  \onslide<3->{
    \textbf{Round 3}
    \begin{itemize}
      \item % TODO: third clue. Players refine and trade.
    \end{itemize}
  }

  \onslide<4->{
    \vspace{0.3cm}
    \begin{sqt:result}
      \textbf{True value:} TODO. Liquidate all positions here.
    \end{sqt:result}
  }
\end{frame}


% #####################################################################
\section{Pricing Futures}
% #####################################################################

\begin{frame}{What is a Future?}
  % TODO: fill in --- agreement now to exchange an asset at a fixed
  % price at a fixed future date. Contrast with a spot trade.
\end{frame}

\begin{frame}{Pricing a Future}
  % TODO: fill in --- fair future price = spot compounded to expiry.
  % F = S * (1 + r)^t. Intuition via no-arbitrage / cost of carry.
\end{frame}

\begin{frame}{Trading Sim 3: Pricing a Basic Future}
  \begin{sqt:definition}
    The asset is a \textbf{future}: it settles at a spot value $S$ compounded
    at a per-period rate $r$ over $t$ periods.
    $$F = S \,(1+r)^{t}.$$
  \end{sqt:definition}

  \begin{itemize}
    \item % TODO: give S, r, t. e.g. S = 100, r = 5\% per period, t = 3.
    \item Make a market on $F$ --- this is really an exercise in estimating
      compound interest in your head.
  \end{itemize}

  \vspace{0.3cm}
  \textcolor{sqt-question}{Estimate $(1+r)^t$ quickly. When is the linear
  approximation $1 + rt$ good enough?}
\end{frame}


% #####################################################################
\section{Pricing Options}
% #####################################################################

\begin{frame}{What is an Option?}
  % TODO: fill in --- the RIGHT but not the obligation to buy (call)
  % at a strike K. Payoff = max(V - K, 0). Why it is never negative.
\end{frame}

\begin{frame}{Trading Sim 4: Option on a d6, Strike 3}
  \begin{sqt:definition}
    A single die is rolled, value $d \sim \mathrm{Unif}\{1,\dots,6\}$.
    The asset is a \textbf{call option} with strike $K = 3$:
    $$\text{payoff} = \max(d - 3,\; 0).$$
  \end{sqt:definition}

  \begin{center}
    \diehidden

    \vspace{0.4cm}
    \begin{tabular}{c|cccccc}
      $d$ & 1 & 2 & 3 & 4 & 5 & 6 \\ \hline
      $\max(d-3,0)$ & 0 & 0 & 0 & 1 & 2 & 3 \\
    \end{tabular}
  \end{center}

  \vspace{0.3cm}
  \textcolor{sqt-question}{What is the fair price of this option?}
  % Fair value = (0+0+0+1+2+3)/6 = 1.
\end{frame}

\begin{frame}{Trading Sim 5: Option on Sum of Three d6, Strike 11}
  \begin{sqt:definition}
    Three dice are rolled, $S = d_1 + d_2 + d_3$.
    The asset is a \textbf{call option} with strike $K = 11$:
    $$\text{payoff} = \max(S - 11,\; 0).$$
  \end{sqt:definition}

  \begin{center}
    \die{4}\hspace{0.6cm}\die{6}\hspace{0.6cm}\die{3}
    % TODO: set / hide dice as appropriate for the sim.
  \end{center}

  \vspace{0.3cm}
  \textcolor{sqt-question}{Fair price? $\mathbb{E}[\max(S-11,0)]$ over the
  distribution of the sum of three dice.}
  % E[S] = 10.5; need the payoff distribution above the strike.
\end{frame}


% #####################################################################
\section{Wrap Up}
% #####################################################################

\begin{frame}{Wrap Up}
  % TODO: fill in --- recap of fair value, spreads, futures, options,
  % and takeaways for trading interviews.
\end{frame}

\end{document}
